% ============================================
% DOC-05 (EN)
% English translation — SPEC-935-R355
% ============================================
\section*{\textbf{DOC-05} --- Session Report II}
\markboth{DOC-05 --- Session Report II}{}

\noindent \textit{Dr. Heitor Oliveira --- March 21, 1979}

\noindent \textbf{Dr. Heitor Oliveira --- March 21, 1979}

\vspace{0.5em}\noindent\rule{\textwidth}{0.5pt}\vspace{0.5em}

\noindent \textbf{Patient:} Joaquim Antônio Correia --- Patient 1,260
\textbf{Type:} Psychiatric assessment session
\textbf{Duration:} 90 minutes (extended)

\vspace{0.5em}\noindent\rule{\textwidth}{0.5pt}\vspace{0.5em}

\noindent \textbf{Patient Observations:}

\noindent The patient presented differently today. More animated. Almost euphoric. He smiled when I entered the room --- an asymmetrical smile, pulled toward the right side of his face, where the yellow eye glows more intensely than in the previous session.

\noindent He asked: \textit{\textquotedblleft{}Did you dream?\textquotedblright{}}

\noindent I answered yes, that I had an unusual dream. He did not ask for details. He only nodded, as if he already knew.

\noindent In this session, Joaquim spoke continuously for 75 minutes. Practically without interruption. He described his childhood in Ceará, the drought of 1915, the walk from Senador Pompeu toward Fortaleza, his father's death, and his confinement at the Alagadiço Camp. The accounts are vivid --- there are sensory details (smells, temperatures, textures) that are unusual even in patients with preserved memory.

\noindent The tone of the account changed when he described the \textquotedblleft{}encounter\textquotedblright{}.

\noindent According to Joaquim, he found a creature in a trench of bodies behind the concentration camp. The description contained psychotic elements (yellow eyes, bone spines, external ribs), but the \textit{internal consistency} is remarkable. He described the encounter with the same conviction with which he described his father's death.

\noindent \textbf{Verbatim record:}

\noindent \textit{\textquotedblleft{}She said: 'You are hungry too.' I said yes. She said: 'I eat for you. I eat for all who cannot eat.' And I sat down. I sat on the edge of the trench and watched her eat. I felt no fear. I felt hunger. A hunger that never went away.\textquotedblright{}}

\noindent \textbf{Analysis:}

\noindent - The \textquotedblleft{}creature\textquotedblright{} appears to be a manifestation of Joaquim's childhood trauma, possibly a personification of the hunger and death he witnessed in childhood.
- The encounter at the trench may represent a dissociative moment: Joaquim may have witnessed something unbearable (the trench of bodies, the death of others) and his mind created an alternative narrative to process the event.
- The \textquotedblleft{}hunger that never goes away\textquotedblright{} is consistent with Post-Traumatic Stress Disorder, in which the traumatic event is relived indefinitely.

\noindent \textbf{Therapeutic Plan:}

\noindent - I requested access to the patient's diary, which was initially denied, but the patient agreed to hand it over \textquotedblleft{}when it is ready\textquotedblright{}.
- Maintain weekly sessions.

\noindent \textbf{Personal Note:}

\noindent I have begun having recurring dreams. They are not ordinary nightmares --- they are extremely vivid, with smells and tactile sensations. I wake up with the taste of dirt in my mouth. I have not mentioned this to Joaquim, but he knows. He knows.

\vspace{0.5em}\noindent\rule{\textwidth}{0.5pt}\vspace{0.5em}

\noindent \textit{Dr. Heitor Oliveira}
\textit{Psychiatrist --- Maximum Security Ward}
\textit{Barbacena Hospital-Colony}

\vspace{0.5em}
\noindent\textit{\textbf{↪ Routes:} \rota{MEM-15} \rotasep \rota{DOC-07} \rotasep \rota{LUC-03}}
